%Nicht anfassen, so ist der Dokumentenaufbau
\documentclass[a4paper, 12pt]{article}
\input{settings/packages}
\input{settings/page}
\begin{document}
\input{settings/titlepage}
\setcounter{page}{1}
%%%%%%%%%%%%%%%%%%%%%%%%%%%%%
%%%Ab hier Inhalt einfügen%%%
%%%%%%%%%%%%%%%%%%%%%%%%%%%%%

\section{Biographisches \& Verortung in Werken} \label{einleitung}
\todo{Unbedingt betonen, dass Whitehead Mathematiker war!} %solche Verweise werden in rot gedruckt und sollten in der Endversion nicht mehr enthalten sein!

\blindtext[3]
%

\section{Hauptteil}

\Blindtext


%

\subsection{Erste Unter"uberschrift}
\blindtext
\subsubsection{Unterunter"uberschrift}

\Blindtext

\subsection{Zweite Unter"uberschrift}

%mit \nocite{} werden Eintr"age dem Literaturverzeichnis hinzugef"ugt, ohne dass ein Verweis auftaucht.
\nocite{Jonas.2010}
\nocite{Whitehead.1987}
\nocite{Jonas.2003}

%zum zitieren in einer Fussnote benutzt man
\footcite[vorgestellter Text (bspw "`vgl."')][nachgestellter Text (bspw. Seitenzahl)]{Whitehead.1987}

Zum Zitieren im Flie"stext eignet sich \textcite[vgl.][S. 333]{Whitehead.1987}. Dadurch wird nur der Name des Autors im Text gedruckt, zusammen mit einer erl"auternden Fu"snote.
\subsection{Im Vergleich}
%
Es ist auch m"oglich, Blockzitate einzuf"ugen. Am sch"onsten wird das so:
\begin{footcitedquote}[\cite{Jonas.2010}]
    \blindtext
\end{footcitedquote}
\section{Fazit}
%

%%%%%%%%%%%%%%%%%%%%%%%%%%%%
%%%%%Ende des Dokuments%%%%%
%%%%%%%%%%%%%%%%%%%%%%%%%%%%
\newpage
\input{settings/endofdocument}
\end{document}